\documentclass[../../../include/open-logic-section]{subfiles}

\begin{document}

\olfileid{sth}{ord-arithmetic}{mult}

\olsection{ক্রমসংখ্যার গুণ}

এবার ক্রমসংখ্যার গুণের আলোচনা করি। এর পদ্ধতি যোগের
অনুরূপ। ধরি, $\alpha$-কে $\beta$ দিয়ে গুণ করতে চাই।
প্রস্থ $\alpha$ ও উচ্চতা $\beta$—এমন একটি আয়তাকার
ঘরবিন্যাস কল্পনা করুন। এক সারি বরাবর এগিয়ে, তার
পরে পরের সারিতে গিয়ে এভাবে পুরো বিন্যাসটি পার হলে
যে ক্রম মেলে, সেটিই $\alpha$ ও $\beta$-র গুণফল।
অন্যভাবে বললে, $\beta$-র \emph{প্রত্যেক} উপাদানের
জায়গায় $\alpha$-র একটি অনুলিপি বসিয়ে
$\alpha$ ও $\beta$-র গুণফল পাওয়া যায়।

এটি কঠোরভাবে প্রকাশ করতে $\alpha$ ও $\beta$-র
কার্তেসীয় গুণফলের উপর বিপরীত অভিধানিক ক্রম ব্যবহার করি:

\begin{defn}
যেকোনো ক্রমসংখ্যা $\alpha, \beta$-র গুণফল
$\alpha \ordtimes \beta = \ordtype{\alpha \times \beta,
\rlexless}$।
\end{defn}
\noindent
আবারও সংজ্ঞাটি সুগঠিত কি না নিশ্চিত করতে হবে:

\begin{lem}\ollabel{ordtimeslessiswo}
যেকোনো ক্রমসংখ্যা $\alpha$ ও $\beta$-র জন্য
$\tuple{\alpha \times \beta, \rlexless}$
একটি সুক্রমবিন্যাস।
\end{lem}

\begin{proof}
\olref[add]{ordsumlessiswo}-এর মতোই প্রমাণ।
\end{proof}
\noindent
গুণের নিচের পুনরাবৃত্তিমূলক ধর্মগুলিও প্রমাণ করা কঠিন নয়:

\begin{lem}\ollabel{ordtimesrecursion}
যেকোনো ক্রমসংখ্যা $\alpha, \beta$-র জন্য:
\begin{align*}
	\alpha \ordtimes 0 &= 0\\
	\alpha \ordtimes (\beta \ordplus 1) &=
		(\alpha \ordtimes \beta) \ordplus \alpha\\
	\alpha \ordtimes \beta &=
		\bigcup_{\delta < \beta}(\alpha \ordtimes \delta) &&
		\text{যখন $\beta$ সীমা ক্রমসংখ্যা}.
\end{align*}
% BN-SRC-804 TARGET-CORRECTION-BEGIN
\emph{উৎস-সংশোধন-নোট।} সীমার দফায় সাধারণ ঊর্ধ্বসীমা,
অর্থাৎ পর্যায়ক্রমিক মানগুলির সংযোগ চাই। মূলের কঠোর
ক্ষুদ্রতম ঊর্ধ্বসীমা শূন্য গুণনীয়কে ভুল: সব আগের মান
শূন্য হলে তার strict ঊর্ধ্বসীমা $1$, অথচ শূন্য ক্ষেত্রের
ক্রমপ্রকার $0$। সংশোধিত সংযোগ $0$-ই দেয়।
$\alpha\neq0$ হলে মানগুলি সহশেষী এবং সর্বাধিক সদস্যহীন,
তখন কঠোর ও সাধারণ ঊর্ধ্বসীমা এখানে মেলে।
% BN-SRC-804 TARGET-CORRECTION-END
\end{lem}

\begin{proof}
অনুশীলনের জন্য রাখা হলো।
\end{proof}

বস্তুত যোগের মতো এখানেও সরাসরি সংজ্ঞা না দিয়ে
এই পুনরাবৃত্তি-সমীকরণগুলি দিয়েই ক্রমসংখ্যার গুণ
সংজ্ঞায়িত করতে পারতাম। যোগের মতোই গুণের কয়েকটি
আচরণও পরিচিত:

\begin{lem}\ollabel{ordinalmultiplicationisnice}
$\alpha, \beta, \gamma$ ক্রমসংখ্যা হলে:
\begin{enumerate}
	\item\ollabel{ordtimes1} $\alpha \neq 0$ এবং
	$\beta < \gamma$ হলে $\alpha \ordtimes \beta
	< \alpha \ordtimes \gamma$;
	\item\ollabel{ordtimes2} $\alpha \neq 0$ এবং
	$\alpha \ordtimes \beta = \alpha\ordtimes\gamma$
	হলে $\beta = \gamma$;
	\item\ollabel{ordtimes3} $\alpha \ordtimes (\beta \ordtimes
	\gamma) = (\alpha \ordtimes \beta) \ordtimes \gamma$;
	\item\ollabel{ordtimes4} $\alpha \leq \beta$ হলে
	$\alpha \ordtimes \gamma \leq \beta \ordtimes \gamma$;
	\item\ollabel{ordtimes5} $\alpha \ordtimes (\beta \ordplus
	\gamma) = (\alpha \ordtimes \beta)\ordplus
	(\alpha\ordtimes \gamma)$।
\end{enumerate}
\end{lem}

\begin{proof}
অনুশীলনের জন্য রাখা হলো।
\end{proof}

ইচ্ছেমতো আরও ফল প্রমাণ করতে বা খুঁজে দেখতে পারেন।
তবে \olref[ord-arithmetic][add]{ordsumnotcommute}
দেখে নিচের ফলে অবাক হওয়ার কথা নয়:

\begin{prop}
ক্রমসংখ্যার গুণ বিনিময়ী নয়:
$2 \ordtimes \omega = \omega < \omega \ordtimes 2$।
\end{prop}

\begin{proof}
$2 \ordtimes \omega = \supstrict_{n < \omega}
(2\ordtimes n) = \omega \in
\supstrict_{n < \omega}(\omega \ordplus n)
= \omega \ordplus \omega = \omega \ordtimes 2$।
\end{proof}
\noindent
এখানেও স্বজ্ঞাগত কারণটি সরল। $2 \ordtimes \omega$
হিসাব করতে প্রত্যেক স্বাভাবিক সংখ্যার জায়গায় দুটি
বস্তু বসান। স্বাভাবিক সংখ্যাগুলির ফাঁকে সব ``অর্ধেক''
সংখ্যা ঢুকিয়ে দিলে, অর্থাৎ
$\Setabs{\nicefrac{n}{2}}{n \in \omega}$-র
স্বাভাবিক ক্রম নিলেও একই ক্রমপ্রকার পাওয়া যাবে।
এভাবে দেখলে ক্রমপ্রকারটি যে $\omega$-রই সমান,
তা স্পষ্ট। কিন্তু $\omega \ordtimes 2$ হিসাব করতে
$\omega$-র দুটি অনুলিপি একটির পরে আরেকটি বসাতে হয়।

\begin{prob}
\olref[sth][ord-arithmetic][mult]{ordtimeslessiswo},
\olref[sth][ord-arithmetic][mult]{ordtimesrecursion}
এবং
\olref[sth][ord-arithmetic][mult]{ordinalmultiplicationisnice}
প্রমাণ করুন।
\end{prob}

\end{document}
